\documentclass[11pt]{ctexart}
\usepackage[a4paper,margin=1in]{geometry}
\usepackage{graphicx}
\usepackage{xcolor}
\usepackage{hyperref}
\definecolor{refgray}{gray}{0.45}
\title{\textbf{Market Views｜全球投行叙事汇编}\\[0.4em]{\large 资金从极端集中向更广泛领域轮动，但基础尚不稳固；亚洲利率与外汇市场结构性分化，AI与半导体需求结构加速演变}}
\author{KC桌面 自动生成}
\date{260722}
\begin{document}
\maketitle
\tableofcontents
\newpage
\section{一页摘要}
\begin{itemize}
  \item 亚洲利率与外汇市场结构性分化：香港长端利率受美债曲线趋平与本地贷款需求疲弱双重压制，存在下行空间；人民币中间价逆周期因子调节幅度收窄，市场预期趋于稳定；G10外汇持仓中挪威克朗与美元头寸反向调整，反映资金从主流货币向北欧小币种轮动。
  \item AI对半导体需求的结构性影响深化：开源模型推高推理算力与HBM需求，混合键合从逻辑芯片向HBM渗透，中国WFE进口回暖但结构分化，存储长期协议尚难消除周期波动。
  \item 网络安全预算增速持续超越IT预算且差距扩大，软件/SaaS为主要受益方向；AI代理技术正在重塑安全治理架构，身份与云安全仍是优先领域，但端点安全增长预期放缓。
  \item 全球消费与零售市场分化加剧：美妆、高端手表及AI驱动的电商平台表现强劲，而大众消费品、运动品牌库存清理及部分区域奢侈品需求仍面临压力。
  \item GLP-1需求驱动多肽CDMO并购升温，三星生物收购PolyPeptide标志亚洲CDMO全球化进入新阶段；美国专业药房以3\%处方量贡献超50\%药费，增速9\%远超传统药房2.5\%。
  \item 地缘政治引发的能源供应中断重塑石油和LNG市场：波斯湾原油流量降至战前水平45\%以下，卡塔尔LNG出口几乎停滞；工业领域呈现结构性分化，数据中心储能需求爆发式增长。
  \item 德国国防预算刚性扩张，以债务融资加速支出，目标2029年占GDP 3.6\%；亚太房地产市场分化，中国新房与二手房成交边际改善但库存高企，开发商盈利分化加剧。
  \item 全球资金正从半导体和大型科技股向更广泛行业扩散，但轮动尚未结束；日本科技板块估值分化至历史极端水平；中国宏观环境呈现更极端的K型分化，政策宽松窗口打开但面临多重约束。
\end{itemize}
\section{全球投行叙事汇编}
今日纳入 55 篇投行/券商研究，按 8 个市场、资产与产业主题系统整理。
\newpage
\subsection{宏观与利率：亚洲市场分化与资金流向新信号}
\textbf{亚洲利率与外汇市场正经历结构性分化：香港长端利率受美债曲线趋平与本地贷款需求疲弱双重压制，存在下行空间；人民币中间价模型显示逆周期因子调节幅度收窄，市场预期趋于稳定；G10外汇持仓中挪威克朗与美元头寸反向调整，反映资金从主流货币向北欧小币种轮动。}
\subsubsection*{主流共识}
\begin{itemize}
  \item 美债曲线可能趋平，若科技股调整持续，长端利率仍有下行空间，进而传导至香港利率市场。
  \item 港元贷款增长未见起色，存贷比持续下降，银行体系资金需求不足对长端利率形成压制。
  \item 人民币中间价模型预测连续低于前值，逆周期因子调节幅度收窄，市场预期趋于稳定。
\end{itemize}
\subsubsection*{分歧与条件差异}
\begin{itemize}
  \item 香港长端利率下行空间：NOM认为5y5y IRS有下行空间（目标3.55\%），但需关注财政储备企稳与债券供给放缓的抵消作用。
  \item G10外汇持仓方向分化：期权市场显示挪威克朗多头与美元空头增加，而期货市场英镑多头与澳元空头增加，不同市场信号存在分歧。
  \item 特步核心品牌与索康尼增速分化：核心品牌销售放缓且库存上升，而索康尼线下仍保持20\%+增长，线上策略保守导致整体增速回落。
\end{itemize}
\subsubsection*{机构观点}
\paragraph{NOM} 香港利率市场长端存在下行空间，建议接收5年期远期5年期港元IRS。逻辑基于三点：美债曲线可能趋平，若科技股调整持续，长端利率仍可能走低；港元贷款增长未见起色，存贷比从73.2\%降至71.0\%，资金需求不足；财政储备同比转正，非金融企业债券发行放缓，金融债发行可能通过互换压低IRS曲线。当前5y5y IRS报3.83\%，目标3.55\%。
{\small\color{refgray}数据：港元存贷比从73.2\%降至71.0\%\par}
{\small\color{refgray}数据：按揭贷款增速从1.0\%升至3.3\%\par}
{\small\color{refgray}数据：港美利差从-85bp回升至-50bp以上\par}
{\small\color{refgray}数据：5y5y港币IRS当前3.83\%，目标3.55\%\par}
{\small\color{refgray}边际变化：此前市场关注美债曲线陡化与港美利差扩大，现转向美债曲线趋平与港美利差收窄，叠加本地贷款需求持续疲弱，长端利率下行逻辑增强。\par}
\paragraph{NOM} 人民币中间价模型预测显示，不含逆周期因子预测值为6.7708，较前日6.7948低240bp，较即期收盘价高38bp；含逆周期因子后升至6.7808，仍较前日实际中间价低140bp。逆周期因子调节幅度收窄，反映市场预期趋于稳定。模型误差持续收窄，预测精度提升。
{\small\color{refgray}数据：不含逆周期因子预测值6.7708，较前值低240bp\par}
{\small\color{refgray}数据：含逆周期因子预测值6.7808，较前值低140bp\par}
{\small\color{refgray}数据：较前日官方即期收盘价高38bp\par}
{\small\color{refgray}数据：美元指数变动是隔夜市场最大贡献因子\par}
{\small\color{refgray}边际变化：此前逆周期因子调节幅度较大，现收窄至140bp，表明央行对市场波动的容忍度上升，中间价定价更贴近市场均衡水平。\par}
\paragraph{NOM} 特步核心品牌2Q26销售额同比中个位数下降，1H26整体低个位数增长，主因需求变化、购物中心渠道不足、竞争加剧及天气因素。库存周转天数从4.5个月升至4.5-5.0个月。索康尼品牌2Q26增速从1Q26的20\%+回落至低个位数，线上策略保守，线下仍保持20\%+增长。
{\small\color{refgray}数据：核心品牌2Q26销售额同比中个位数下降\par}
{\small\color{refgray}数据：1H26整体销售额同比低个位数增长\par}
{\small\color{refgray}数据：库存周转天数从4.5个月升至4.5-5.0个月\par}
{\small\color{refgray}数据：索康尼2Q26增速从20\%+回落至低个位数\par}
{\small\color{refgray}边际变化：此前市场关注特步核心品牌增长与索康尼高增速，现核心品牌销售放缓且库存上升，索康尼线上策略保守导致整体增速回落，增长动能减弱。\par}
\paragraph{MS} 截至7月17日当周，G10外汇持仓出现分化：期权市场显示投资者增加挪威克朗兑欧元净多头与美元指数净空头，战术性资金持有NOK多头与EUR空头配对；期货市场增加英镑多头与澳元空头。加元情绪改善最为明显。美元投机仓位小幅下降至51.3\%，整体仍偏多。
{\small\color{refgray}数据：期权市场：挪威克朗兑欧元净多头增加\par}
{\small\color{refgray}数据：期权市场：美元指数净空头增加\par}
{\small\color{refgray}数据：期货市场：英镑多头增加\par}
{\small\color{refgray}数据：期货市场：澳元空头增加\par}
{\small\color{refgray}边际变化：此前市场关注美元多头与欧元空头，现转向挪威克朗多头与美元空头，且期权与期货市场信号分化，反映资金从主流货币向北欧小币种轮动。\par}
\subsubsection*{关键数据}
\begin{itemize}
  \item 港元存贷比从73.2\%降至71.0\%
  \item 港美利差从-85bp回升至-50bp以上
  \item 5y5y港币IRS当前3.83\%，目标3.55\%
  \item 人民币中间价不含逆周期因子预测值6.7708，较前值低240bp
  \item 逆周期因子调节幅度收窄至140bp
  \item 特步核心品牌2Q26销售额同比中个位数下降
  \item 索康尼2Q26增速从20\%+回落至低个位数
  \item 美元投机仓位降至51.3\%
\end{itemize}
\begin{figure}[htbp]
\centering
\includegraphics[width=0.88\linewidth]{figures/fig_045.jpg}
\caption{Figure 1 - 1. Potential US curve flattening and HK-US spread no longer too wide. Over the past week, the market's focus has been the renewed rise in oil prices and the correction in some technology stocks. We believe if concerns over the correction in equity prices persi}
\end{figure}
\begin{figure}[htbp]
\centering
\includegraphics[width=0.88\linewidth]{figures/fig_141.jpg}
\caption{Exhibit 1 - Exhibit 1: MS \& CO. LLC Molly Nickolin Strategist}
\end{figure}
\begin{figure}[htbp]
\centering
\includegraphics[width=0.88\linewidth]{figures/fig_046.jpg}
\caption{Figure 2 - 3. Fiscal reserves have stabilized, while non-financial HK bond issuance has slowed. While we remain cautious on Hong Kong's fiscal situation over medium term, because projects such as Northern Metropolis could require a significant amount of investment over m}
\end{figure}
\begin{figure}[htbp]
\centering
\includegraphics[width=0.88\linewidth]{figures/fig_142.jpg}
\caption{Exhibit 2 - Exhibit 3: Exhibit 4:}
\end{figure}
{\small\color{refgray}本节综合 4 篇研究；涉及 NOM、MS。}
\newpage
\subsection{AI与半导体：需求结构分化，技术路径加速}
\textbf{AI对半导体需求的结构性影响正在深化：开源模型推高推理算力与HBM需求，混合键合从逻辑芯片向HBM渗透，中国WFE进口回暖但结构分化，存储长期协议（LTA）尚难消除周期波动。}
\subsubsection*{主流共识}
\begin{itemize}
  \item 开源大模型（如Kimi K3）的普及增加而非减少对HBM和计算资源的总需求，英伟达等算力供应商持续受益。
  \item 混合键合在逻辑芯片封装（如台积电SoIC）中已进入规模化部署，2028年相关设备市场将显著增长。
  \item AI驱动北美企业生产率提升约9.6\%，但净岗位流失约5\%，行业分化明显（银行/硬件受益，软件/专业服务承压）。
\end{itemize}
\subsubsection*{分歧与条件差异}
\begin{itemize}
  \item HBM用混合键合的时间表：Bernstein认为三星可能在2027年HBM4E中率先采用，但若仅12层堆叠则推迟至2028年HBM5；而UBS强调开源模型的长上下文窗口将加速HBM需求，但未明确技术节点。
  \item 存储长期协议（LTA）能否平滑周期：Bernstein测算当前担保仅覆盖未来收入的0.6\%，认为无法消除盈利波动；而部分市场观点可能高估LTA的保护力度。
  \item 中国WFE进口前景：Bernstein指出上半年累计同比下降10\%，但6月光刻设备环比增长196\%，下半年存储设备进口预计走强；而ASML预计中国区收入占比将从33\%降至20\%，显示供应端约束仍存。
  \item 苹果AI服务器芯片性能：MS认为自研芯片（M2 Ultra/Baltra ASIC）面临能效与性能的权衡，可能影响Siri AI体验或增加外部云依赖；而苹果供应链信息显示SolC产能未取消，但量产时间在2027年后。
\end{itemize}
\subsubsection*{机构观点}
\paragraph{Bernstein} 混合键合设备市场2028年将增长至约14亿美元，逻辑芯片是主要驱动力：台积电SoIC产能预计从2025年1万片/月增至2028年8.5万片/月，每千片/月需5-6台键合设备。HBM领域采用时间表存在不确定性，三星可能在HBM4E（16层）中率先使用1层混合键合，若仅12层则推迟至HBM5（2028年），HBM每千片/月仅需2-3台设备。
{\small\color{refgray}数据：2028年混合键合设备市场14亿美元（逻辑8.87亿+HBM 4.89亿）\par}
{\small\color{refgray}数据：台积电SoIC产能2028年达8.5万片/月（2025年1万片/月）\par}
{\small\color{refgray}数据：HBM每千片/月产能仅需2-3台键合设备（逻辑需5-6台）\par}
{\small\color{refgray}边际变化：此前市场对混合键合关注集中于逻辑芯片，Bernstein首次系统量化HBM渗透路径及设备需求差异，并上调2028年总市场预测。\par}
\paragraph{Bernstein} 基于韩国海关出口数据，三星二季度HBM收入可能达67亿美元，环比增长89\%，较此前预测高出25\%。6月忠清南道出口环比飙升93\%，首次单月超过SK海力士。SK海力士二季度出口环比增长19\%，符合预期，价值/重量比持平，暂未看到HBM4出货信号。HBM定价与常规存储价格出现分化，2027年LTA谈判预计将带动盈利预期上调。
{\small\color{refgray}数据：三星2Q26 HBM收入预估67亿美元（环比+89\%，较预测高25\%）\par}
{\small\color{refgray}数据：6月忠清南道出口环比+93\%，首次单月超SK海力士\par}
{\small\color{refgray}数据：SK海力士2Q26出口环比+19\%，价值/重量比持平\par}
{\small\color{refgray}边际变化：此前市场对三星HBM增长预期偏保守，出口数据显示HBM4已开始放量，且月度季节性后置（55\%出口集中在6月），需调整线性外推框架。\par}
\paragraph{Bernstein} 2026年6月中国WFE进口34亿美元，环比+57\%，同比+1\%，略高于2025年月均32亿美元。上半年累计进口160亿美元，同比-10\%，主要受光刻（-18\%）和干法刻蚀（-14\%）拖累。6月光刻设备进口8.42亿美元，环比+196\%，已回升至2025年月均水平。下半年存储设备进口预计明显走强，区域供应格局变化：美国+新加坡+马来西亚份额从2024年33\%升至2026年YTD 43\%，日本从26\%降至21\%。
{\small\color{refgray}数据：6月WFE进口34亿美元（环比+57\%，同比+1\%）\par}
{\small\color{refgray}数据：上半年累计160亿美元（同比-10\%）\par}
{\small\color{refgray}数据：6月光刻设备进口8.42亿美元（环比+196\%）\par}
{\small\color{refgray}数据：美国+新加坡+马来西亚进口份额升至43\%（2024年33\%）\par}
{\small\color{refgray}边际变化：此前市场担忧中国WFE进口持续疲软，6月数据表明光刻设备已从4月低点恢复，且区域供应链转移趋势加速。\par}
\paragraph{Bernstein} 新内存长期协议（LTA）能提升盈利底部，但无法消除周期波动。当前披露担保总额约330亿美元，而同期内存供应商年收入预期约1.3万亿美元，覆盖3-5年收入需保护约5.2万亿美元，担保仅占0.6\%。即使乐观情景，LTA也只能将低谷期营业利润率从历史最差的-10\%\textasciitilde{}10\%提升至0\%\textasciitilde{}30\%。约30-50\%的市场份额（消费电子、交易型买家）难以被LTA覆盖。
{\small\color{refgray}数据：当前LTA担保总额330亿美元，仅覆盖未来收入的0.6\%\par}
{\small\color{refgray}数据：乐观情景下LTA将低谷期营业利润率从-10\%\textasciitilde{}10\%提升至0\%\textasciitilde{}30\%\par}
{\small\color{refgray}数据：30-50\%的市场份额难以被LTA覆盖\par}
{\small\color{refgray}边际变化：此前市场部分观点认为LTA将根本性改变内存周期，Bernstein通过担保/收入比例测算指出保护力度有限，盈利波动性仍将存在。\par}
\paragraph{UBS} 开源模型Kimi K3（2.8万亿参数）发布，强调规模效应而非低价竞争，其长上下文窗口和始终在线推理模式将增加对HBM和存储的需求。存储子行业到2028年累计自由现金流占市值比例最高，美光尤为突出。市场对MU和SK海力士的估值差已缩小至约0.3倍EV/S，但MU在DRAM密度、LP-DDR和能效方面技术领先未被充分定价。
{\small\color{refgray}数据：Kimi K3参数2.8万亿，100万token上下文窗口\par}
{\small\color{refgray}数据：存储子行业到2028年累计FCF/市值比例最高\par}
{\small\color{refgray}数据：MU与SK海力士EV/S估值差缩小至0.3倍\par}
{\small\color{refgray}边际变化：此前市场担忧开源模型降低算力需求，UBS指出开源模型更长的上下文窗口和广泛部署反而推高HBM需求，且存储FCF创造能力与市场定价背离。\par}
\paragraph{MS} 苹果自研AI服务器芯片（M2 Ultra/Baltra ASIC）面临性能挑战，其能效优先的设计理念与AI推理所需的高功耗、高吞吐存在结构性差异。首款AI专用芯片Baltra（与博通合作）瞄准2027上半年小规模量产，属试验性产品；第二代AI ASIC预计2028年量产，功耗包络可能接近商用AI GPU。若自研芯片性能不足，Siri AI复杂查询表现将受影响，或促使苹果将更多工作负载迁移至Google Cloud，增加成本。
{\small\color{refgray}数据：Baltra ASIC瞄准2027上半年小规模量产\par}
{\small\color{refgray}数据：第二代AI ASIC预计2028年量产，功耗接近商用AI GPU\par}
{\small\color{refgray}数据：苹果SolC产能未取消，但同时服务于M5 Max/Ultra和AI服务器\par}
{\small\color{refgray}边际变化：此前市场对苹果AI芯片进展较为乐观，MS指出其技术积累与AI推理需求存在结构性差异，且量产时间表晚于Siri AI功能上线（2026年秋季），形成时间差约束。\par}
\paragraph{MS} 2026年AlphaWise调查显示，AI在过去12个月驱动平均5\%净岗位流失和9.6\%生产率提升。行业分化显著：银行和科技硬件/半导体是净受益者（AI提升运营效率并创造增量需求），软件生产率提升最明显（10.4\%）但净岗位流失最高（7\%），专业服务类似。2-5年经验的中层员工受影响最大，应届生反而较小。
{\small\color{refgray}数据：AI驱动净岗位流失5\%（12\%裁员+15\%不补缺-22\%新招聘）\par}
{\small\color{refgray}数据：生产率提升9.6\%\par}
{\small\color{refgray}数据：软件行业净岗位流失7\%（最高），生产率提升10.4\%（最高）\par}
{\small\color{refgray}数据：2-5年经验中层员工受影响最大\par}
{\small\color{refgray}边际变化：此前市场对AI就业影响多为定性讨论，MS通过808家企业调查提供定量数据，并揭示行业结构性分化，指出银行实际岗位削减可能高于调查显示的4\%。\par}
\paragraph{MS} 微软Azure和Copilot即将迎来增速拐点。上调Azure FY28/FY29收入预期（高于共识5\%/7.8\%），基于新建货币化模型：Azure AI的Rev/MW在保守情景下990万美元（毛利率25\%），乐观情景下2277万美元（毛利率35\%），市场目前按基础资源供应商定价，忽略了更高层AI服务拉动的收入乘数。Copilot变现路径从单一席位转向“三引擎”：席位+E7升级+消费型AI服务，正从试点进入企业级大规模部署。
{\small\color{refgray}数据：Azure FY28/FY29收入预期高于共识5\%/7.8\%\par}
{\small\color{refgray}数据：Azure AI Rev/MW：保守990万美元，乐观2277万美元\par}
{\small\color{refgray}数据：Copilot变现三引擎：席位+E7升级+消费型AI服务\par}
{\small\color{refgray}边际变化：此前市场关注微软AI的容量约束，MS指出拐点已至，且市场低估了Azure AI的收入乘数和Copilot的ARPU扩张潜力。\par}
\paragraph{NOM} Kimi K3（2.8万亿参数开源模型）在Artificial Analysis智能指数排名第三，每任务成本0.94美元，约为Opus 4.8的一半、Fable 5的三分之一，但较前代K2.6上涨4倍，标志月之暗面从低价转向高性价比定位。中国模型在全球开发者token用量占比已从一年前不到2\%升至超过45\%。主权AI趋势下，中美头部大模型厂商均有增长空间，竞争加剧将拉动先进计算和网络需求。
{\small\color{refgray}数据：Kimi K3参数2.8万亿，智能指数排名第三\par}
{\small\color{refgray}数据：每任务成本0.94美元（Opus 4.8的一半，Fable 5的三分之一）\par}
{\small\color{refgray}数据：中国模型全球token用量占比超45\%（一年前不到2\%）\par}
{\small\color{refgray}边际变化：此前市场关注中国模型的价格战，K3表明竞争转向高性价比，且中国模型已从跟随者变为全球开发者生态主要供给方，对芯片采购和模型定价产生结构性影响。\par}
\subsubsection*{关键数据}
\begin{itemize}
  \item 2028年混合键合设备市场14亿美元（逻辑8.87亿+HBM 4.89亿）
  \item 三星2Q26 HBM收入预估67亿美元（环比+89\%，较预测高25\%）
  \item 6月中国WFE进口34亿美元（环比+57\%），上半年累计同比-10\%
  \item LTA担保仅覆盖未来收入的0.6\%，乐观情景下低谷期利润率从-10\%\textasciitilde{}10\%提升至0\%\textasciitilde{}30\%
  \item AI驱动净岗位流失5\%，生产率提升9.6\%
  \item Azure AI Rev/MW：保守990万美元，乐观2277万美元
  \item 中国模型全球token用量占比超45\%（一年前不到2\%）
  \item 苹果Baltra ASIC瞄准2027上半年小规模量产
\end{itemize}
\begin{figure}[htbp]
\centering
\includegraphics[width=0.88\linewidth]{figures/fig_011.jpg}
\caption{EXHIBIT 1 - EXHIBIT 2: Our hybrid bonder estimates are between Besi's mid and high case guidance. \#\# LOGIC: CLEAR TRAJECTORY OF HYBRID ADOPTION Hybrid bonding is expanding in adoption among the various interconnect technologies, especially f}
\end{figure}
\begin{figure}[htbp]
\centering
\includegraphics[width=0.88\linewidth]{figures/fig_019.jpg}
\caption{Exhibit 1 - EXHIBIT 1: Being a close proxy for HBM shipment, Korea multichip memory export to Taiwan and Malaysia reached a historical high at USD 6.2B in Jun 2026. EXHIBIT 2: Export from S. Chungcheong was up 93\% MoM and 75\% QoQ vs. Mar. Mu}
\end{figure}
\begin{figure}[htbp]
\centering
\includegraphics[width=0.88\linewidth]{figures/fig_039.jpg}
\caption{Exhibit 1 - EXHIBIT 1: June import was \textbackslash{}\$3.4 billion Total WFE imports to China by month (USDmn) ■ Jun-25 ■ Jun-26 ● YoY (RHS) EXHIBIT 2: Month-over-month showed 57\% increase, largely driven by lithography import. Total WFE imports to China}
\end{figure}
\begin{figure}[htbp]
\centering
\includegraphics[width=0.88\linewidth]{figures/fig_084.jpg}
\caption{EXHIBIT 5 - EXHIBIT 5: Current LTA Deposit Levels Appear Insufficient to Structurally Reduce Earnings Volatility}
\end{figure}
{\small\color{refgray}本节综合 9 篇研究；涉及 Bernstein、UBS、MS、NOM。}
\newpage
\subsection{网络安全预算加速扩张，AI治理成为新控制点}
\textbf{网络安全支出增速持续超越IT预算，且差距扩大，软件/SaaS为主要受益方向；AI代理技术正在重塑安全治理架构，身份与云安全仍是优先领域，但端点安全增长预期放缓。}
\subsubsection*{主流共识}
\begin{itemize}
  \item 网络安全预算增速显著快于IT预算，且这一趋势在2026年中期调查中较2025年底进一步强化。
  \item 软件/SaaS是网络安全支出增长的最大受益方，70\%的CISO预期未来12个月软件支出将增长。
  \item 身份安全、云安全和端点安全是研究优先级最高的三个领域。
\end{itemize}
\subsubsection*{分歧与条件差异}
\begin{itemize}
  \item AI对安全支出的影响：部分观点认为AI带来净新增预算，另一些则担心模型幻觉和准确率问题可能限制投入。
  \item 端点安全增长预期出现下降，而SecOps、身份安全和威胁情报的增量支出意愿提升明显。
  \item SSE（安全服务边缘）支出增长预期最弱，且较上次调查进一步走弱，与云安全持续高增长形成对比。
\end{itemize}
\subsubsection*{机构观点}
\paragraph{Bernstein} 2026年6月对100位美国CISO的半年度调查显示，43\%的受访者预计网络安全支出增速快于IT支出，仅20\%认为IT更快，差距较2025年底调查扩大。超过半数CISO上调全年支出预期，暗示下半年支出将明显强于上半年。金融、科技和服务业年化增速约6\%，所有行业均超4\%。网络安全支出占IT预算比例维持在8\%左右，但行业间差异收窄。软件/SaaS是最大受益方，70\%预期增长；身份安全、云安全和端点安全优先级最高，但端点安全未来12个月支出计划增幅下降，SSE增速预期最弱。
{\small\color{refgray}数据：43\%的CISO预计网络安全支出增速快于IT，仅20\%认为IT更快\par}
{\small\color{refgray}数据：超过半数CISO上调全年支出预期\par}
{\small\color{refgray}数据：金融、科技和服务业年化增速约6\%\par}
{\small\color{refgray}数据：网络安全支出占IT预算比例维持在8\%\par}
{\small\color{refgray}边际变化：与2025年底调查相比，网络安全支出增速领先IT的差距进一步扩大，且更多CISO上调全年预期，表明下半年支出可能集中释放。\par}
\paragraph{Bernstein} 2026年6月西欧经历有记录以来最热六月，连续热浪将区域供冷基础设施推至极限。巴黎拉德芳斯商务区冰储罐在周末储备100\%容量，但一天内消耗近三分之二，夜间无法补足；供水温度从4.5°C升至7.5°C。全球区域供冷市场2025年312亿美元，预计2034年达615亿美元，年复合增速8\%。中东非洲占42\%份额，欧洲仅7\%，渗透率偏低。法国目标2030年供冷量弹性提升至2TWh以上，2035年达2.5-3.0TWh。极端天气正在测试现有网络承载上限，同时为低碳冷却方案提供加速扩张窗口。
{\small\color{refgray}数据：2026年6月西欧为有记录以来最热六月\par}
{\small\color{refgray}数据：巴黎拉德芳斯冰储罐一天内消耗近三分之二\par}
{\small\color{refgray}数据：供水温度从4.5°C升至7.5°C\par}
{\small\color{refgray}数据：全球区域供冷市场2025年312亿美元，2034年预计615亿美元，CAGR 8\%\par}
{\small\color{refgray}边际变化：极端天气频发正在加速区域供冷决策，欧洲追赶空间最大，法国2030年目标成为关键锚点。\par}
\paragraph{JPM} Intapp与SS\&C正从不同起点向代理型AI转型，但均试图成为受监管客户部署AI时必须经过的治理层。Intapp从云原生垂直SaaS出发，将Celeste定位为受治理的代理型平台，核心差异化在于“上下文+治理+工作流嵌入”，结合企业专有机构记忆与信息隔离墙控制。SS\&C则从大规模智能自动化和业务流程外包切入，以WorkHQ作为Blue Prism之上的编排层，依托企业级护栏与合规优先治理。治理不是一项功能，而是采用瓶颈，两家公司都在争夺这一控制点。
{\small\color{refgray}数据：Intapp Celeste平台强调模型无关的代理型工作流\par}
{\small\color{refgray}数据：SS\&C拥有2.3万客户，从内部验证AI降本增效向客户推广\par}
{\small\color{refgray}数据：Intapp差异化在于上下文、治理与工作流嵌入的组合\par}
{\small\color{refgray}数据：SS\&C依托企业级防护和合规优先架构\par}
{\small\color{refgray}边际变化：AI代理技术从通用工具转向行业特定治理层，Intapp和SS\&C分别从SaaS和BPO出发，争夺受监管客户AI部署的控制点。\par}
\subsubsection*{关键数据}
\begin{itemize}
  \item 43\%的CISO预计网络安全支出增速快于IT，仅20\%认为IT更快
  \item 超过半数CISO上调全年支出预期
  \item 70\%的CISO预期未来12个月软件支出将增长
  \item 网络安全支出占IT预算比例维持在8\%
  \item 端点安全支出计划增幅下降，为降幅第二大类别
  \item 全球区域供冷市场2025年312亿美元，2034年预计615亿美元，CAGR 8\%
  \item 欧洲仅占全球区域供冷市场7\%份额
  \item Intapp Celeste平台强调模型无关的代理型工作流
\end{itemize}
\begin{figure}[htbp]
\centering
\includegraphics[width=0.88\linewidth]{figures/fig_077.jpg}
\caption{Exhibit 17 - EXHIBIT 1: Our mid-year 2026 survey shows a little less extreme variation in cybersecurity spend as a \% of IT than the EOY 2025 survey did. The overall average remained about the same. Total annual cyber spending as \% of IT EXHIB}
\end{figure}
\begin{figure}[htbp]
\centering
\includegraphics[width=0.88\linewidth]{figures/fig_080.jpg}
\caption{EXHIBIT 4 - EXHIBIT 5: In mid-2026, a greater portion of CISOs anticipate cybersecurity outgrowing IT than they did even six months ago at our EOY 2025 survey. Which budget is growing faster? \% of respondents}
\end{figure}
\begin{figure}[htbp]
\centering
\includegraphics[width=0.88\linewidth]{figures/fig_006.jpg}
\caption{EXHIBIT 1 - EXHIBIT 1: June 2026 was the warmest June on record for western Europe (anomalies in surface air temperature for June 2026) EXHIBIT 2: Daily surface air temperature for western Europe \#\# HOW DOES A DISTRICT COOLING SYSTEM (DCS)}
\end{figure}
\begin{figure}[htbp]
\centering
\includegraphics[width=0.88\linewidth]{figures/fig_007.jpg}
\caption{EXHIBIT 1 - EXHIBIT 3: General scheme of a district cooling system}
\end{figure}
{\small\color{refgray}本节综合 3 篇研究；涉及 Bernstein、JPM。}
\newpage
\subsection{消费与零售：分化加剧，结构性机遇浮现}
\textbf{全球消费与零售市场呈现显著分化：美妆、高端手表及AI驱动的电商平台表现强劲，而大众消费品、运动品牌库存清理及部分区域奢侈品需求仍面临压力。企业正通过定价策略、渠道优化和技术创新应对挑战，行业竞争从规模扩张转向效率与品牌价值提升。}
\subsubsection*{主流共识}
\begin{itemize}
  \item 美妆品类成为消费亮点，增速显著高于家庭护理等大众消费品，消费者支出向特定品牌和产品集中。
  \item 东南亚电商平台（Shopee、TikTok Shop）正从补贴竞争转向运营效率，通过削减补贴和调整配送费实现价格上行。
  \item 瑞士手表出口回暖，美国市场是主要驱动力，大中华区需求分化（内地调整、香港受益于境外消费）。
\end{itemize}
\subsubsection*{分歧与条件差异}
\begin{itemize}
  \item 茅台飞天提价后批价迅速反弹，显示核心产品定价权稳固；但非标产品发货量同比下降约30\%，定价能力仍在修复中。
  \item 运动品牌Puma库存清理进度不一：英国DTC渠道SKU数量下降近50\%，而美国市场仍处于回购阶段，进度落后。
  \item AI购物助手对电商格局的影响存在分歧：MS认为大型平台将受益，而市场担忧其颠覆现有格局。
\end{itemize}
\subsubsection*{机构观点}
\paragraph{Bernstein} 达美乐Q2订单量增长显著，但客单价承压，美国同店销售仅增0.1\%。下半年新品成败是关键变量，但历史测试与市场反应常有偏差。国际业务受DPE清理促销影响，同店下降0.1\%。当前约18倍预期市盈率已反映悲观预期，但不确定性仍存。
{\small\color{refgray}数据：美国同店销售Q2仅增0.1\%\par}
{\small\color{refgray}数据：订单量增长显著，外卖和外带均正增长\par}
{\small\color{refgray}数据：国际同店下降0.1\%\par}
{\small\color{refgray}数据：美国净开店目标从175+下调至约150\par}
{\small\color{refgray}边际变化：市场关注点从订单量增长转向客单价和加盟商盈利能力，新品成败成为下半年核心变量。\par}
\paragraph{Bernstein} 瑞士手表6月出口额同比增11.2\%，但剔除工作日效应后仅增1.1\%。美国市场增12.7\%是主要驱动力，大中华区内地降16.5\%，香港增6.9\%，消费向境外转移。中东增23.8\%部分源于补货，可持续性存疑。200-500瑞郎区间量价齐升，中档表（500-3000瑞郎）收缩。
{\small\color{refgray}数据：6月瑞士手表出口额24亿瑞郎，同比增11.2\%\par}
{\small\color{refgray}数据：日历调整后实际增速1.1\%\par}
{\small\color{refgray}数据：美国市场增12.7\%\par}
{\small\color{refgray}数据：中国内地降16.5\%，香港增6.9\%\par}
{\small\color{refgray}边际变化：全球手表需求逐步调整，美国市场企稳，大中华区消费向境外转移，中档表需求承压。\par}
\paragraph{Bernstein} Puma库存清理在英国取得实质进展，DTC渠道SKU数量较峰值下降近50\%，折扣策略为深度大于广度。美国市场仍处回购阶段，SKU数量未降，但批发合作伙伴如Footlocker的SKU已大幅下降（降80\%）。下一阶段关键从清库存转向品牌复兴与产品创新。
{\small\color{refgray}数据：英国DTC SKU数量较峰值下降近50\%\par}
{\small\color{refgray}数据：Footlocker SKU下降80\%\par}
{\small\color{refgray}数据：Mostro系列超80\% SKU打折\par}
{\small\color{refgray}边际变化：库存清理重点从英国转向美国，品牌复兴与产品创新成为下一阶段核心。\par}
\paragraph{Bernstein} 东南亚电商价格竞争趋于缓和。Shopee通过削减补贴实现价格上行（印尼最终价环比涨6\%），TikTok Shop减少补贴但降低配送费缓冲（印尼最终价环比涨10\%）。两家平台策略同步调整，从争夺订单转向保护自身经济模型。
{\small\color{refgray}数据：Shopee印尼最终价环比涨6\%\par}
{\small\color{refgray}数据：TikTok Shop印尼最终价环比涨10\%\par}
{\small\color{refgray}数据：Shopee补贴折扣贡献-11\%降幅\par}
{\small\color{refgray}数据：TikTok Shop补贴折扣从-3\%收窄\par}
{\small\color{refgray}边际变化：行业从补贴换增长转向关注可持续运营，平台策略罕见同步调整。\par}
\paragraph{Bernstein} 美国家庭与个人护理（HPC）行业分化，美妆是亮点。宝洁美妆部门Q3至今销售增6.4\%（销量增4.7\%），远超整体2.5\%增速。高露洁整体增3.5\%，宠物业务支撑。e.l.f. Beauty护肤品类增38.8\%，Naturium品牌销售增58.2\%。家庭护理品类表现平稳。
{\small\color{refgray}数据：宝洁美妆Q3至今销售增6.4\%\par}
{\small\color{refgray}数据：宝洁整体销售增2.5\%\par}
{\small\color{refgray}数据：高露洁整体增3.5\%\par}
{\small\color{refgray}数据：e.l.f. Beauty护肤品类增38.8\%\par}
{\small\color{refgray}边际变化：美妆品类韧性持续，消费者支出向特定品牌集中，家庭护理品类相对平稳。\par}
\paragraph{GS} 茅台第五轮市场化改革，飞天出厂价/零售价上调7.9\%/6.5\%。提价后批价迅速反弹55元至1705元/瓶。二季度终端动销稳定，经销商库存仅约2周，全年发货进度超去年同期（65-70\% vs 55\%）。非标产品发货量同比下降约30\%，定价权仍在修复。
{\small\color{refgray}数据：飞天出厂价/零售价上调7.9\%/6.5\%\par}
{\small\color{refgray}数据：原箱飞天批价涨55元至1705元/瓶\par}
{\small\color{refgray}数据：经销商库存约2周\par}
{\small\color{refgray}数据：全年发货进度65-70\%（去年同期55\%）\par}
{\small\color{refgray}边际变化：茅台转向更频繁、市场化的定价操作，飞天定价权稳固，非标产品仍在控量保价。\par}
\paragraph{GS} 瑞士手表6月出口额同比增11.2\%至24亿瑞郎，美国市场增13\%是最大驱动力，中国内地降17\%（5月降21\%），香港增7\%。历峰集团大中华区恢复增长，斯沃琪上半年大中华区增9\%。200-500瑞郎区间表现最佳（量增51\%），中档表（500-3000瑞郎）下滑。
{\small\color{refgray}数据：6月瑞士手表出口额24亿瑞郎，同比增11.2\%\par}
{\small\color{refgray}数据：美国市场增13\%\par}
{\small\color{refgray}数据：中国内地降17\%\par}
{\small\color{refgray}数据：香港增7\%\par}
{\small\color{refgray}边际变化：中国消费者手表购买力向境外转移，美国市场在关税冲击后逐步恢复。\par}
\paragraph{GS} Ulta促销活动客流与网页流量显示营销策略有效。Big Summer Beauty Sale客流增10.9\%，Spring Haul客流增2.1\%，所有活动网页流量均录得双位数增长。Jumbo Love与21 Days of Beauty客流分别下降1.7\%和1.8\%。市场关注Q2最难的同比基数，但数据表明营销投入正带动实际到店行为。
{\small\color{refgray}数据：Big Summer Beauty Sale客流增10.9\%\par}
{\small\color{refgray}数据：Spring Haul客流增2.1\%\par}
{\small\color{refgray}数据：Jumbo Love客流降1.7\%\par}
{\small\color{refgray}数据：21 Days of Beauty客流降1.8\%\par}
{\small\color{refgray}边际变化：Ulta营销投入有效性得到第三方数据验证，客流和网页流量增长表明策略正在见效。\par}
\paragraph{MS} AI驱动的购物助手（Agentic Commerce）将成为电商下一阶段演进动力。到2030年线上零售渗透率将从22\%升至27\%，对应约7万亿美元GMV，年复合增长率从7\%提升至9\%。Agentic Commerce有望贡献约6\%增量市场空间。拥有规模、用户关系和生态的平台（如Amazon、Walmart、Alibaba）处于更有利位置。
{\small\color{refgray}数据：2030年线上零售渗透率预计从22\%升至27\%\par}
{\small\color{refgray}数据：2030年全球电商GMV约7万亿美元\par}
{\small\color{refgray}数据：年复合增长率从7\%提升至9\%\par}
{\small\color{refgray}数据：Agentic Commerce贡献约6\%增量\par}
{\small\color{refgray}边际变化：AI购物助手从概念走向实际应用，大型电商平台有望受益，而非被颠覆。\par}
\subsubsection*{关键数据}
\begin{itemize}
  \item 达美乐美国同店销售Q2仅增0.1\%
  \item 瑞士手表6月出口额24亿瑞郎，同比增11.2\%
  \item Puma英国DTC SKU数量较峰值下降近50\%
  \item Shopee印尼最终价环比涨6\%
  \item 宝洁美妆Q3至今销售增6.4\%
  \item 飞天批价涨55元至1705元/瓶
  \item Ulta Big Summer Beauty Sale客流增10.9\%
  \item 2030年全球电商GMV预计达7万亿美元
\end{itemize}
\begin{figure}[htbp]
\centering
\includegraphics[width=0.88\linewidth]{figures/fig_001.jpg}
\caption{EXHIBIT 1 - EXHIBIT 1: DPZ Same-store Sales Growth DPZ Same-store Sales Growth EXHIBIT 2: DPZ Domestic Unit Growth EXHIBIT 3: DPZ Same-store Sales Growth}
\end{figure}
\begin{figure}[htbp]
\centering
\includegraphics[width=0.88\linewidth]{figures/fig_016.jpg}
\caption{EXHIBIT 1 - EXHIBIT 1: Exports of Swiss Wristwatches watches by price categories Swiss watch exports by price - Total market EXHIBIT 2: Swiss watch exports by material Swiss watch exports by material - Total market (CHF Mil) EXHIBIT 3: Swi}
\end{figure}
\begin{figure}[htbp]
\centering
\includegraphics[width=0.88\linewidth]{figures/fig_052.jpg}
\caption{EXHIBIT 1 - EXHIBIT 1: Puma UK - Markdown Breadth and Markdon Depth \% Puma UK - Markdown Breadth and Markdown Depth (\%) EXHIBIT 2: Puma UK - SKU count EXHIBIT 3: Puma UK - Average Price (£)}
\end{figure}
\begin{figure}[htbp]
\centering
\includegraphics[width=0.88\linewidth]{figures/fig_057.jpg}
\caption{EXHIBIT 1 - EXHIBIT 1: Mix of products cheaper in a platform ID: Indonesia; PH: Phillipines; VN: Vietnam; MY: Malaysia; TH: Thailand EXHIBIT 2: April- June 2026: Change in final product prices was in the -5\% to 30\% range for Sea compared to}
\end{figure}
{\small\color{refgray}本节综合 9 篇研究；涉及 Bernstein、GS、MS。}
\newpage
\subsection{医疗与制药：GLP-1外包扩张与美国专业药房结构分化}
\textbf{GLP-1需求驱动多肽CDMO并购升温，三星生物收购PolyPeptide标志亚洲CDMO全球化进入新阶段；美国专业药房以3\%处方量贡献超50\%药费，增速9\%远超传统药房2.5\%，但集中度高、政策风险低；联合健康MA业务通过主动收缩会员实现利润率超预期回升，MLR降至86.7\%，但Optum Insights和商业险改善尚未被市场充分定价。}
\subsubsection*{主流共识}
\begin{itemize}
  \item GLP-1类药物全球需求持续高增长，多肽CDMO成为战略资产，并购是获取技术know-how和产能的高效路径。
  \item 美国专业药房年支出增速约9\%，远高于传统药房的2.5\%，主要由新药上市、处方量增加和单位成本上升共同推动。
  \item 联合健康MA业务通过定价纪律和会员结构优化实现利润率修复，2026年全年MA利润率有望回升至3\%以上。
\end{itemize}
\subsubsection*{分歧与条件差异}
\begin{itemize}
  \item 三星生物收购PolyPeptide对国内CDMO的短期影响：NOM认为短期内不会构成实质性冲击，因PolyPeptide增速（10\%）远低于药明康德TIDES（49\%），且整合需要时间。
  \item 专业药房与PBM的业务边界：Bernstein强调专业药房不参与PBM返利安排，政策风险较低；但若监管推动PBM剥离旗下专业药房，市场格局可能重塑。
  \item 联合健康商业险和Medicaid业务的恢复节奏：Bernstein认为2026年商业险恢复只是延迟而非倒退，全面回归要到2027年后；Medicaid仍在调整中，利润率低于-1\%。
\end{itemize}
\subsubsection*{机构观点}
\paragraph{NOM} 三星生物以2.7万亿韩元收购PolyPeptide，对应2026年EV/EBITDA 19.1倍，估值与龙沙、Bachem接近（18.6倍），战略上具有吸引力。PolyPeptide参与全球1/3多肽三期临床，1H26营收2.366亿欧元（+41.6\%），EBITDA利润率20.7\%。收购旨在获取多肽CDMO技术know-how，短期内不会对药明康德TIDES（2025年收入约110亿元，3年CAGR 49\%）构成直接冲击。
{\small\color{refgray}数据：收购价2.7万亿韩元，对应EV/EBITDA 19.1倍\par}
{\small\color{refgray}数据：PolyPeptide 1H26营收2.366亿欧元，同比+41.6\%\par}
{\small\color{refgray}数据：EBITDA利润率20.7\%，2026全年营收增速指引25\%-30\%\par}
{\small\color{refgray}数据：药明康德TIDES 2025年收入约110亿元，3年CAGR 49\%\par}
{\small\color{refgray}边际变化：三星生物从抗体CDMO跨界进入多肽CDMO，标志GLP-1外包需求推动亚洲CDMO全球化进入新阶段；此前市场预期三星生物将自建产能，收购路径显示多肽CDMO进入壁垒高于预期。\par}
\paragraph{Bernstein} 美国专业药房仅占处方量3\%-4\%，却贡献超50\%药品支出，年增速9\% vs 传统药房2.5\%。专业药房不参与PBM返利安排，政策风险较低；规模效应和有限分销网络构成主要竞争壁垒，CVS、Accredo、Optum三家主导市场。四种运营模型（传统、罕见病、家庭输注、医生诊所给药）各有不同服务与定价模式。
{\small\color{refgray}数据：专业药房占处方量3\%-4\%，贡献超50\%药品支出\par}
{\small\color{refgray}数据：年支出增速9\%，传统药房仅2.5\%\par}
{\small\color{refgray}数据：CVS、Accredo、Optum三家占据主导地位\par}
{\small\color{refgray}边际变化：此前市场更多关注PBM监管风险，Bernstein首次系统区分专业药房与PBM的业务边界，指出专业药房对监管变化的敏感度相对较低，且增长驱动力来自结构性因素（新药、处方量、成本上升）而非一次性事件。\par}
\paragraph{Bernstein} 联合健康2Q26调整后EPS 6.38美元，超市场共识30.4\%；MLR 86.7\%，比预期低170bp，同比降270bp。MA业务通过主动减少不盈利会员（预计全年减少110万人）实现利润率回升至3\%以上。商业险利润率目前4\%（目标7\%），Medicaid仍在调整中（利润率低于-1\%），Optum Insights AI工具已在70\%医生中落地，目标年底90\%。
{\small\color{refgray}数据：2Q26调整后EPS 6.38美元，超共识30.4\%\par}
{\small\color{refgray}数据：MLR 86.7\%，比预期低170bp，同比降270bp\par}
{\small\color{refgray}数据：2026全年MA会员预计减少110万人，利润率回升至3\%以上\par}
{\small\color{refgray}数据：商业险利润率4\%，目标7\%；Medicaid利润率低于-1\%\par}
{\small\color{refgray}边际变化：市场此前预期MA利润率修复需更长时间，但2Q26数据表明定价纪律和会员优化效果快于预期；同时Optum Insights的AI商业化潜力（从内部工具转向外部客户）是未被充分定价的新变量。\par}
\subsubsection*{关键数据}
\begin{itemize}
  \item 三星生物收购PolyPeptide对应EV/EBITDA 19.1倍，与龙沙、Bachem的18.6倍接近
  \item PolyPeptide 1H26营收2.366亿欧元，同比+41.6\%，EBITDA利润率20.7\%
  \item 药明康德TIDES 2025年收入约110亿元，3年CAGR 49\%
  \item 专业药房占处方量3\%-4\%，贡献超50\%药品支出，年增速9\% vs 传统2.5\%
  \item 联合健康2Q26 MLR 86.7\%，比预期低170bp，同比降270bp
  \item 联合健康2026全年MA会员预计减少110万人，利润率回升至3\%以上
  \item 商业险利润率4\%，目标7\%；Medicaid利润率低于-1\%
  \item Optum Health AI工具覆盖70\%医生，目标年底90\%
\end{itemize}
\begin{figure}[htbp]
\centering
\includegraphics[width=0.88\linewidth]{figures/fig_072.jpg}
\caption{EXHIBIT 2 - EXHIBIT 2: What is specialty pharmacy - growth Specialty cost trend of 9\% vs 2.5\% for traditional over recent period Drivers of growth – mix, price inflation and increased number of prescriptions At an industry level, about 1/3}
\end{figure}
\begin{figure}[htbp]
\centering
\includegraphics[width=0.88\linewidth]{figures/fig_087.jpg}
\caption{EXHIBIT 1 - EXHIBIT 1: UNH - Adj. EPS UHC - Operating margin, \% EXHIBIT 2: UHC - Operating margins}
\end{figure}
\begin{figure}[htbp]
\centering
\includegraphics[width=0.88\linewidth]{figures/fig_089.jpg}
\caption{EXHIBIT 3 - EXHIBIT 3: UNH Relative NTM PE \#\# 2Q 2026 RESULTS EXHIBIT 4: UnitedHealth 2Q 2026 Result vs. Bernstein and Consensus Estimates}
\end{figure}
\begin{figure}[htbp]
\centering
\includegraphics[width=0.88\linewidth]{figures/fig_073.jpg}
\caption{EXHIBIT 2 - EXHIBIT 2: What is specialty pharmacy - growth Specialty cost trend of 9\% vs 2.5\% for traditional over recent period Drivers of growth – mix, price inflation and increased number of prescriptions At an industry level, about 1/3}
\end{figure}
{\small\color{refgray}本节综合 4 篇研究；涉及 NOM、Bernstein。}
\newpage
\subsection{能源与工业：供应冲击与结构性分化}
\textbf{地缘政治引发的能源供应中断正在重塑石油和LNG市场，同时工业领域呈现显著的结构性分化：数据中心储能需求爆发式增长，而传统工业如电梯、汽车和玻纤则面临需求疲软或结构性调整。}
\subsubsection*{主流共识}
\begin{itemize}
  \item 地缘政治风险导致能源供应不确定性显著上升，波斯湾原油流量降至战前水平45\%以下，卡塔尔LNG出口几乎停滞。
  \item 数据中心成为电池储能需求增长最快的下游场景，美国数据中心电力容量预计到2030年将增长至103GW。
  \item 中国市场需求疲软，对全球工业品价格和能源需求形成压制。
\end{itemize}
\subsubsection*{分歧与条件差异}
\begin{itemize}
  \item 对油价上行空间的判断存在分歧：GS认为市场对极端情景的定价弹性已高于战争初期，上行空间有限；而Bernstein和JPM则强调供应中断可能持续，风险偏向上行。
  \item 对卡塔尔LNG冬季供应恢复的预期不同：JPM认为市场定价过于乐观，卡塔尔可能在2-3周内被迫停产；而部分市场参与者仍假设冬季前能恢复。
  \item 工业领域分化明显：数据中心储能和电子纱需求强劲，而传统工业（电梯、汽车、粗纱）面临需求下行压力。
\end{itemize}
\subsubsection*{机构观点}
\paragraph{Bernstein} 美国数据中心正成为电池储能需求增长最快的下游场景，平均配置时长已升至约5小时，在建项目中24\%采用离网模式。数据中心电力容量假设从41GW上调至2030年的103GW，累计电池需求达277GWh。年度数据中心电池需求将从当前的5GWh增长至2030年的66GWh，增幅约10倍。新项目的电池容量与数据中心装机容量比值已达4至8倍，显示后续需求仍有上行空间。
{\small\color{refgray}数据：美国数据中心电力容量2030年预计达103GW\par}
{\small\color{refgray}数据：累计电池需求277GWh\par}
{\small\color{refgray}数据：年度电池需求从5GWh增至66GWh\par}
{\small\color{refgray}数据：新项目电池容量比达4-8倍\par}
{\small\color{refgray}边际变化：数据中心储能时长从4小时提升至5-8小时，离网模式占比上升，推动电池需求预期大幅上调。\par}
\paragraph{Bernstein} Schindler二季度有机收入增速仅1.1\%，低于市场预期160个基点，主要受中国市场持续调整拖累。订单增长2.9\%略超预期，现代化改造业务订单增速约11\%是主要驱动力。利润率13.5\%同比提升90个基点，但收入缺口使全年低至中个位数增长指引面临挑战，下半年需实现约5\%增长才能达标。
{\small\color{refgray}数据：Q2有机收入增速1.1\%，低于预期160bps\par}
{\small\color{refgray}数据：订单增长2.9\%，现代化改造订单增长11\%\par}
{\small\color{refgray}数据：利润率13.5\%，同比提升90bps\par}
{\small\color{refgray}数据：EMEA新梯订单H1增长12.8\%\par}
{\small\color{refgray}边际变化：现代化改造业务正在抵消新安装业务的结构性下滑，但收入端持续失速使全年指引可信度下降。\par}
\paragraph{Bernstein} 沃尔沃汽车二季度营收低于预期6\%，EBIT低于预期25\%，并取消全年销量增长指引。业绩下滑主要由产品组合向低价车型（EX30、EX40）倾斜驱动，插电混动销量同比下降12\%。中国市场零售销量下降37\%，欧洲市场持平，美国市场增长9\%。成本削减计划执行好于预期，上半年已完成全年5亿瑞典克朗目标，但不足以抵消外部压力。
{\small\color{refgray}数据：Q2营收777亿瑞典克朗，低于预期6\%\par}
{\small\color{refgray}数据：EBIT 8.26亿克朗，低于预期25\%\par}
{\small\color{refgray}数据：中国零售销量下降37\%\par}
{\small\color{refgray}数据：上半年自由现金流流出52亿克朗\par}
{\small\color{refgray}边际变化：产品组合下移成为利润压缩主因，公司撤回全年销量增长指引，但成本削减领先于计划。\par}
\paragraph{Citi} 中国玻纤市场延续粗纱承压、电子纱偏强的分化格局。粗纱2400tex缠绕直接纱均价3768.5元/吨，周环比下降0.26\%，下游加工企业开工率偏低，库存累积。电子纱G75价格稳定于16800-17500元/吨，电子布8.5-8.7元/米，供需偏紧、低库存和强劲的CCL需求支撑价格。新增产能爬坡较慢，高端产品供应有限，中期仍有上行空间。
{\small\color{refgray}数据：粗纱均价3768.5元/吨，周环比-0.26\%\par}
{\small\color{refgray}数据：G75电子纱16800-17500元/吨\par}
{\small\color{refgray}数据：7628电子布8.5-8.7元/米\par}
{\small\color{refgray}数据：122条产线运行，总产能903.8万吨/年\par}
{\small\color{refgray}边际变化：电子纱供应瓶颈不在于总产能，而在于新增产能爬坡慢和高端产品可用量有限，供需偏紧格局持续。\par}
\paragraph{GS} 波斯湾原油流量已降至战前水平45\%以下，布伦特期货曲线略高于GS对2026Q4和2027年的80/75美元预测。油价不确定性偏向上行，但市场对极端情景的定价弹性已高于战争初期。Q2库存去化超300万桶/日，OECD柴油和SPR库存偏低。中国6月原油净海运进口同比下降44\%，但约20亿桶库存缓冲和煤炭替代能力为价格设置隐形天花板。欧洲柴油时间价差具备三重结构性支撑。
{\small\color{refgray}数据：波斯湾原油流量降至战前水平45\%以下\par}
{\small\color{refgray}数据：Q2库存去化超300万桶/日\par}
{\small\color{refgray}数据：中国6月原油净海运进口同比下降44\%\par}
{\small\color{refgray}数据：中国持有约20亿桶原油库存\par}
{\small\color{refgray}边际变化：市场假设了更大的需求弹性和对低库存的更高容忍度，油价上行空间低于战争初期估算。\par}
\paragraph{JPM} 7月12日后霍尔木兹海峡LNG运输几乎完全停止，卡塔尔装载率从30\%降至22\%。已装载货物大部分成为浮动仓储。卡塔尔在满负荷条件下约有5天岸上储存能力，以当前20\%利用率可延至20-25天。预计2-3周内可能达到储罐满溢状态，届时可能被迫停产，重启需2-3个月，将影响北半球冬季供应。TTF价格约60欧元/兆瓦时，隐含的不确定性溢价仍不充分。
{\small\color{refgray}数据：7月12日后霍尔木兹海峡LNG运输几乎停止\par}
{\small\color{refgray}数据：卡塔尔装载率从30\%降至22\%\par}
{\small\color{refgray}数据：岸上储存能力约5天（满负荷）或20-25天（20\%利用率）\par}
{\small\color{refgray}数据：波斯湾内仍有23艘LNG船\par}
{\small\color{refgray}边际变化：市场对卡塔尔冬季供应的定价可能过于乐观，储罐满溢时间窗口约2-3周，停产风险上升。\par}
\subsubsection*{关键数据}
\begin{itemize}
  \item 美国数据中心电力容量2030年预计达103GW
  \item 数据中心年度电池需求从5GWh增至66GWh
  \item Schindler Q2收入增速1.1\%，低于预期160bps
  \item 沃尔沃Q2 EBIT低于预期25\%，中国零售销量下降37\%
  \item 粗纱均价3768.5元/吨，周环比-0.26\%
  \item G75电子纱16800-17500元/吨
  \item 波斯湾原油流量降至战前水平45\%以下
  \item 中国6月原油净海运进口同比下降44\%
  \item 卡塔尔LNG装载率从30\%降至22\%
  \item TTF价格约60欧元/兆瓦时
\end{itemize}
\begin{figure}[htbp]
\centering
\includegraphics[width=0.88\linewidth]{figures/fig_024.jpg}
\caption{EXHIBIT 2 - EXHIBIT 2: Total capacity is expected to reach 103 GW in 2030, with cumulative battery capacity reaching 277 GWh. EXHIBIT 3: Project-Level Benchmark for US Data Center BESS}
\end{figure}
\begin{figure}[htbp]
\centering
\includegraphics[width=0.88\linewidth]{figures/fig_044.jpg}
\caption{EXHIBIT 1 - EXHIBIT 1: Schindler Results Scorecard}
\end{figure}
\begin{figure}[htbp]
\centering
\includegraphics[width=0.88\linewidth]{figures/fig_094.jpg}
\caption{Figure 1 - Not for distribution in the People's Republic of China, excluding the Hong Kong Special Administrative Region and Qualified Foreign Institutional Investors. Figure 1. Glass fiber price Figure 2. E-fabric price Figure 3. Glass fiber inventory in key regions}
\end{figure}
\begin{figure}[htbp]
\centering
\includegraphics[width=0.88\linewidth]{figures/fig_121.jpg}
\caption{Exhibit 14 - Exhibit 1: Lower Gulf Flows Means Higher Crude Prices \#\# Upside Price Risks Exhibit 2: Oil Price Risks Are Tilted to the Upside}
\end{figure}
{\small\color{refgray}本节综合 6 篇研究；涉及 Bernstein、Citi、GS、JPM。}
\newpage
\subsection{区域市场与房地产：德国国防扩张与亚太地产分化}
\textbf{德国国防预算的刚性扩张与亚太房地产市场的结构性分化是当前两大主线。德国正以债务融资加速国防支出，目标2029年占GDP 3.6\%，而中国房地产市场呈现新房与二手房成交边际改善但库存高企、开发商盈利分化的格局。}
\subsubsection*{主流共识}
\begin{itemize}
  \item 德国国防预算扩张具有政治共识，且通过债务融资，2029年目标明确。
  \item 中国新房与二手房成交量连续两周环比改善，但领先指标显示需求基础不稳。
  \item 中国房地产开发商盈利分化加剧，高端项目销售强劲但历史低毛利项目持续拖累利润。
\end{itemize}
\subsubsection*{分歧与条件差异}
\begin{itemize}
  \item 德国国防采购策略被视为‘次优’，集中于传统武器和本国企业，而非前沿技术。
  \item 中国房地产市场中，二手房韧性明显强于新房，但库存去化周期仍处高位。
  \item 中国金茂销售强劲但利润承压，而龙湖集团开发业务毛利率转负且修复空间有限。
\end{itemize}
\subsubsection*{机构观点}
\paragraph{DB} 德国正以远超预期的速度推进国防预算扩张，2027年国防与安全预算草案已分配1600亿欧元（GDP的2.7\%），较2026年的2.1\%显著上升。到2029年，国防预算计划升至GDP的3.6\%，核心预算中每三欧元就有一欧元用于国防。2025-2029年累计国防支出约7300亿欧元，较去年秋季计划增加近900亿欧元，其中至少5000亿欧元通过新增债务融资。这一扩张受政治共识支撑，且将国防支出从债务刹车中剥离具有制度性影响。
{\small\color{refgray}数据：2027年国防预算1600亿欧元，占GDP 2.7\%\par}
{\small\color{refgray}数据：2029年国防预算占GDP 3.6\%\par}
{\small\color{refgray}数据：2025-2029年累计国防支出约7300亿欧元\par}
{\small\color{refgray}数据：至少5000亿欧元通过新增债务融资\par}
{\small\color{refgray}边际变化：国防预算增幅比去年秋季制定中期财政计划时更为激进，额外追加220亿欧元国防支出全部通过债务融资。\par}
\paragraph{GS} 中国房地产7月第三周新房成交面积环比增长7\%，同比增长5\%；二手房成交面积环比增长1\%，同比增长7\%，连续第二周同步改善。二手房挂牌量环比下降4\%，低于6月均值8\%，业主涨价预期走强。但新房搜索活动环比微降0.2\%，二手房认购量基本持平，表明成交改善更多来自供给端调整而非需求全面回暖。库存去化周期降至26.9个月，仍高于6月均值。
{\small\color{refgray}数据：新房成交面积环比+7\%，同比+5\%\par}
{\small\color{refgray}数据：二手房成交面积环比+1\%，同比+7\%\par}
{\small\color{refgray}数据：二手房挂牌量环比-4\%，低于6月均值8\%\par}
{\small\color{refgray}数据：库存去化周期26.9个月\par}
{\small\color{refgray}边际变化：新房与二手房成交量连续第二周同步改善，但领先指标（新房搜索、二手房认购）走弱，显示改善持续性存疑。\par}
\paragraph{GS} 中国金茂上半年合同销售同比增长8\%至580亿元，优于覆盖国有房企平均持平水平，主要来自一线城市核心高端项目。但土地购置金额仅占合同销售的约20\%，低于2025年全年的约50\%和GS全年预期的35\%。GS上调2026年合同销售预测3\%至1200亿元，但下调2026年营业收入预测2\%，交付毛利率下调约1个百分点至12\%，反映历史低毛利项目持续入账。2026年净利润预计同比下降7\%。
{\small\color{refgray}数据：上半年合同销售同比+8\%至580亿元\par}
{\small\color{refgray}数据：土地购置金额占合同销售约20\%\par}
{\small\color{refgray}数据：2026年合同销售预测上调3\%至1200亿元\par}
{\small\color{refgray}数据：2026年交付毛利率下调至12\%\par}
{\small\color{refgray}边际变化：上半年拿地节奏明显放缓，但全年拿地预算大概率不变，下半年将加速补充土储。\par}
\paragraph{GS} 龙湖集团2026年上半年合同销售额同比下滑53\%至170亿元，远低于覆盖开发商平均15\%的降幅，主因约80\%可售资源为早期获取的老项目且集中于非核心城市。GS将2026年开发业务毛利率预测调整至-16\%，认为短期内修复空间有限。全年合同销售目标下调28\%至260亿元，对应去化率35\%，预计全年净亏损24亿元（此前预测亏损4.71亿元）。土地储备保持静默，补库节奏延后至2027年。
{\small\color{refgray}数据：上半年合同销售额同比-53\%至170亿元\par}
{\small\color{refgray}数据：约80\%可售资源为早期老项目\par}
{\small\color{refgray}数据：2026年开发业务毛利率预测-16\%\par}
{\small\color{refgray}数据：预计全年净亏损24亿元\par}
{\small\color{refgray}边际变化：开发业务毛利率转负且修复空间有限，土地储备保持静默，补库节奏延后。\par}
\paragraph{DB} 2026年第二季度，DB覆盖的八只相关市场工业及自动化标的外资持股比例均值回升至7.5\%，较一季度环比提高0.8个百分点。这是该系列数据自2024年8月披露频率从日度调整为季度以来，首次出现较为明显的整体改善。恒力液压与双环传动是主要拉动力量，分别录得3.7和2.1个百分点的季度环比增长。这份季度跟踪报告的核心信号是：外资对部分中国工业公司的配置意愿正在修复...
{\small\color{refgray}数据：恒力液压的外资持股比例在二季度达到12.8\%，较一季度上升3.7个百分点，已是连续第三个季度环比走高。这一水平已接近2023年中期约14\%的历史高点。报告没有解释具体驱动因素，但从时间序列看，恒力自2024年四季度起便呈现持续改善趋势，与2022至2023年的节奏变化周期形成对比。该公司的外资持股比例在覆盖标的中排名第二，仅次于汇川技术。\par}
{\small\color{refgray}数据：KC评论： 恒力液压的外资持股修复路径较为完整——从低点连续回升三个季度，且已接近历史高位。这与其他多数标的仍远低于历史峰值的情况不同，说明外资对该公司的配置逻辑可能已发生结构性变化，而非短期交易行为。\par}
\paragraph{GS} 2026年世界人工智能大会（WAIC）的展馆面积同比扩大43\%，参展商超过1100家，但真正让产业观察者驻足的是具身智能展区的变化：242家参展商，较上年增长203\%；208个具身智能产品，同比增长247\%；超过300台人形机器人进行现场演示。这些数字本身已经说明，人形机器人正从一个概念验证阶段加速进入产业构建期。GS在连续四年跟踪WAIC后，今年给出的核心...
{\small\color{refgray}数据：“VLA+WAM”的双系统架构已成为行业主流标准。几乎所有具身智能展台都在强调，视觉-语言-行动模型（VLA）负责视觉理解与任务规划，世界行动模型（WAM）处理物理动力学与状态预测，两者结合形成“系统1”快速运动控制与“系统2”慢速认知推理的分层结构。这种架构共识的建立，意味着人形机器人的AI模型路线图正在收敛，不再像前两年那样各自为政。\par}
{\small\color{refgray}数据：但在数据收集策略上，分歧开始显现。非具身化的第一人称视频数据采集方法已取代传统遥操作，成为行业共识。轻量化采集设备大量涌现，包括手持式双指夹爪、第一人称视觉硬件和外骨骼数据手套。全球非具身数据量预计将从2026年的千万小时级，在3至5年内达到亿小时级，目标是触发模型能力的规模定律。然而，部分团队强调，高质量数据——特别是包含精确触觉和力控制反馈的数据——比单...\par}
\paragraph{GS} 这份GS研报对华润万象生活2026年上半年业绩给出了一个清晰的判断：公司有望实现双位数净利润增长，主要动力来自商业板块的稳定扩张和租金增长。报告同时小幅下调了2026至2028年每股盈利预测，平均约1\%，主要反映对写字楼和增值服务收入更保守的预期。
{\small\color{refgray}数据：报告认为，华润万象生活1H26的运营现金流覆盖比率预计维持在0.6至0.7倍，与历史同期水平一致。尽管第三方住宅项目的现金回款可能面临行业性不确定性，但公司整体分红预期不变，仍为核心利润的100\%派息率（含特别股息）。这一判断与公司管理层此前给出的2026年全年双位数增长指引一致。\par}
{\small\color{refgray}数据：研报指出，华润万象生活1H26的商业板块扩张节奏符合预期，正在推进其“十五五”期间100个新购物中心的管理目标。一个关键参照数据是，母公司华润置地1H26的租金收入同比增长了13\%。报告认为，这一增速为华润万象生活的商业运营收入提供了直接支撑。\par}
\paragraph{GS} GS在2026年7月的一份研报中调整了对Urban Outfitters（URBN）的观察框架。核心判断是：市场对Anthropologie品牌执行不确定性的定价已较为充分，而Urban Outfitters品牌的利润修复和新兴增长引擎的规模化，正在改变公司的整体回报结构。这份报告并非简单看好单一品牌，而是认为组合内各业务的不确定性回报已发生位移。
{\small\color{refgray}数据：报告指出，Urban Outfitters品牌在经历了数年仍在调整的同店销售和利润率趋势后，于2026财年出现正向转折。GS的品牌动量追踪数据显示，该品牌的净购买意愿和搜索趋势均在改善。消费者对款式、价格和品质的评价在过去数月持续走高。从财务角度看，若该品牌利润率恢复到历史中个位数水平，到2028财年将带来超出市场一致预期的利润空间。这一修复主要来自降价幅度...\par}
{\small\color{refgray}数据：KC评论： 市场此前对URBN的关注点集中在Anthropologie的波动上，但UO品牌的利润修复可能成为被低估的变量。报告中的品牌动量数据（净购买意愿、搜索趋势）是判断修复是否持续的实时指标，比财务数据更具前瞻性。\par}
\paragraph{JPM} 中国锂贸易在2026年6月呈现一个清晰信号：碳酸锂进口量同比大幅增长，同时进口价格连续第二个月走高。JPM最新发布的6月中国进出口数据显示，当月碳酸锂进口量达到25.9千吨，较上年同期的17.7千吨增长46\%。价格方面，6月碳酸锂进口均价为每吨19,695美元，环比上升4\%，同比则几乎弹性较高——去年同期仅为每吨10,121美元。
{\small\color{refgray}数据：这些数字叠加在一起，指向一个正在发生的结构性变化：中国锂资源的外部供应不仅在量上持续扩大，其采购成本也在系统性抬升。对于关注锂行业供需格局的读者而言，这组数据值得放在更长的时间维度里审视。\par}
{\small\color{refgray}数据：从累计数据看，2026年前六个月中国碳酸锂净进口量达到177千吨，较2025年同期的115.3千吨增长53.5\%。这一增速远超2025年全年2.8\%的净进口增幅，也高于2024年55\%的增速。JPM报告中的历史数据显示，中国碳酸锂净进口量从2017年的29.2千吨一路攀升至2025年的237.7千吨，十年间扩大超过七倍。2026年上半年延续了这一扩张趋势，且...\par}
\subsubsection*{关键数据}
\begin{itemize}
  \item 德国2029年国防预算占GDP 3.6\%
  \item 中国75城新房成交面积环比+7\%，同比+5\%
  \item 中国20城二手房成交面积环比+1\%，同比+7\%
  \item 中国金茂上半年合同销售同比+8\%至580亿元
  \item 龙湖集团上半年合同销售同比-53\%至170亿元
  \item 龙湖2026年开发业务毛利率预测-16\%
  \item 中国碳酸锂6月进口量同比+46\%至25.9千吨
  \item 中国氢氧化锂上半年转为净进口8.9千吨
\end{itemize}
\begin{figure}[htbp]
\centering
\includegraphics[width=0.88\linewidth]{figures/fig_097.jpg}
\caption{Figure 1 - Figure 1: Germany plans to reach a NATO quota of 3.5\% by 2029 at the latest Figure 2: The bulk of additional defence spending will be debt-funded Robin Winkler Chief Economist}
\end{figure}
\begin{figure}[htbp]
\centering
\includegraphics[width=0.88\linewidth]{figures/fig_106.jpg}
\caption{Exhibit 1 - Exhibit 1: Primary GFA sold last week was +7\% wow and +5\% yoy in c.75 cities Exhibit 2: Primary GFA sold YTD on average was -12\% yoy in c.75 cities, and -13\%/-37\% vs. 2024/2023 level Secondary market: Week 29/YTD volumes were \$}
\end{figure}
\begin{figure}[htbp]
\centering
\includegraphics[width=0.88\linewidth]{figures/fig_098.jpg}
\caption{Figure 2 - Robin Winkler Chief Economist Sebastian-B Becker Senior Economist EUR 640bn earmarked in last autumn's financial plan. The bulk of these defence expenditures, at least EUR 500bn, will be directly funded via fresh debt, thanks to the generous exemption from the}
\end{figure}
\begin{figure}[htbp]
\centering
\includegraphics[width=0.88\linewidth]{figures/fig_107.jpg}
\caption{Exhibit 1 - Exhibit 3: Secondary GFA sold last week was \$+1\textbackslash{}\%\$ wow and \$+7\textbackslash{}\%\$ yoy in c.20 cities Average weekly volume of secondary property sales}
\end{figure}
{\small\color{refgray}本节综合 9 篇研究；涉及 DB、GS、JPM。}
\newpage
\subsection{主题与行业综合：资金轮动、政策博弈与结构性分化}
\textbf{全球市场正经历资金从极端集中向更广泛领域的轮动，但轮动基础尚不稳固；中国宏观环境呈现更极端的“K型”分化，政策宽松窗口打开但面临多重约束；美国医疗政策中期选举后可能转向，日本科技板块估值分化至历史极端水平。}
\subsubsection*{主流共识}
\begin{itemize}
  \item 全球资金正从半导体和大型科技股向更广泛的行业扩散，但轮动尚未结束，指数层面震荡整理可能持续。
  \item 日本科技板块上半年表现极端分化，大型股主导涨幅，7月以来出现轮动迹象，但领涨与落后板块间缺乏明确的盈利或估值拐点。
  \item 中国二季度GDP环比增速放缓，财政脉冲转负是主要拖累，政策宽松窗口正在打开，但决策者在安全、科技竞争和财政约束间面临取舍。
\end{itemize}
\subsubsection*{分歧与条件差异}
\begin{itemize}
  \item 关于美元资金结构：DB认为美元正从“地缘政治锚定”转向“科技周期锚定”，外国官方对美债需求持续走弱，而私人股权资本涌入创纪录；GS则关注日本储蓄回流对日债的潜在支撑，但认为实际影响有限，更多是信号效应。
  \item 关于美国市场领导力扩散：MS认为半导体板块修正尚未结束，超大规模云厂商相对价值优于半导体；Bernstein则认为日本科技板块的轮动需要更有选择性地参与，而非全面切换。
  \item 关于中国政策有效性：GS认为市场型、长期型政策更有效，如降低房贷利率至租金回报率以下；而Bernstein关注美国医疗政策，认为民主党若赢得众议院，将优先逆转特朗普时期的Medicaid和Marketplace削减措施。
\end{itemize}
\subsubsection*{机构观点}
\paragraph{Bernstein} 日本科技板块上半年等权重口径上涨14\%，市值加权口径上涨60\%，半导体、工厂自动化和电子设备大幅跑赢，而IT服务、软件和娱乐持续承压。7月以来资金从领涨板块流向落后板块，但轮动需要更有选择性地参与，因为领涨板块估值和情绪已处于极端水平，而落后板块估值虽有吸引力，但两者间尚未出现明确的盈利或估值拐点。板块整体市盈率18.2倍，低于十年均值0.2个标准差，但市销率2.5倍，较五年均值高出2.6个标准差。盈利上调周期接近峰值，但相对市场仍有空间。
{\small\color{refgray}数据：日本科技板块1H26等权重+14\%，市值加权+60\%\par}
{\small\color{refgray}数据：板块市盈率18.2倍，低于十年均值0.2个标准差\par}
{\small\color{refgray}数据：市销率2.5倍，较五年均值高2.6个标准差\par}
{\small\color{refgray}数据：估值最贵的五分之一公司市销率溢价从十年均值167\%收窄至86\%\par}
{\small\color{refgray}边际变化：7月以来出现轮动迹象，但Bernstein认为轮动基础不稳固，需精选方向而非全面切换。\par}
\paragraph{Bernstein} Seven \& i Holdings一周内完成两笔方向相反的交易：从软银牵头的财团获得至多3000亿日元用于防御收购，同时出资数百亿日元购入波兰最大便利店连锁Żabka的两位数股权。Bernstein认为这两笔交易是必需品而非可选项——日本便利店客流量Q1同比下降0.9\%，提价掩盖了门店空置，销售管理费用膨胀打破了加盟模式的运营杠杆。北美业务尚未完成转型，管理层又增加了欧洲运营负担。公司价值不在于增长，而在于控制权本身的可出售性，上行依赖伊藤忠商事等潜在买家回归。
{\small\color{refgray}数据：Seven \& i从软银财团获至多3000亿日元（约19亿美元）\par}
{\small\color{refgray}数据：日本便利店Q1客流量同比-0.9\%\par}
{\small\color{refgray}数据：Żabka拥有约2.2万家门店\par}
{\small\color{refgray}数据：自由流通股从不到50\%升至52\%\par}
{\small\color{refgray}边际变化：一周内完成防御性融资和海外扩张两笔方向相反的交易，管理层需向股东证明2万亿日元回购计划与资本流出之间的平衡。\par}
\paragraph{Bernstein} 美国软饮市场Q3至今（截至7月11日）品类增长3.3\%，较Q2的4.0\%和Q1的6.0\%放缓，但仍高于消费必需品板块普遍预期。头部公司分化扩大：可口可乐Q3至今销售增长5.2\%，较Q2的4.2\%加速，软饮销售增长7.4\%，市场份额同比提升1.5个百分点至38.8\%，Fairlife销售增长24.2\%，单价为每盎司0.17美元，是软饮品类的近三倍；百事可乐同期销售下滑1.1\%，从Q2的持平转为负增长，饮料和零食业务均承压。
{\small\color{refgray}数据：美国软饮市场Q3至今品类增长3.3\%（Q2为4.0\%，Q1为6.0\%）\par}
{\small\color{refgray}数据：可口可乐Q3至今销售+5.2\%（Q2为+4.2\%），其中销量+3.0\%，价格/组合+2.1\%\par}
{\small\color{refgray}数据：可口可乐软饮市场份额同比+1.5pct至38.8\%\par}
{\small\color{refgray}数据：Fairlife销售+24.2\%，单价0.17美元/盎司，为软饮品类0.06美元的近三倍\par}
{\small\color{refgray}边际变化：品类增速放缓但头部公司分化加剧，可口可乐加速而百事转负，反映产品组合、渠道能力和品牌资产的结构性差异。\par}
\paragraph{Bernstein} 美国中期选举临近，Bernstein认为民主党夺回众议院控制权概率较高，参议院归属仍悬而未决。若民主党至少拿下众议院，医疗政策将出现方向性调整，优先序包括：逆转特朗普时期的Medicaid和Marketplace削减措施、扩大保险覆盖、改善可负担性，以及AI与垂直整合等新兴监管议题。Medicaid工作要求条款是民主党重点攻击目标，已有25个州加哥伦比亚特区联合起诉特朗普政府。CBO数据显示OBBBA法案预计2025-2034年间削减约8860亿美元Medicaid支出，影响约750万人失去保险。
{\small\color{refgray}数据：民主党夺回众议院概率较高，参议院归属悬而未决\par}
{\small\color{refgray}数据：OBBBA法案预计2025-2034年削减约8860亿美元Medicaid支出\par}
{\small\color{refgray}数据：约750万人可能因Medicaid削减失去保险\par}
{\small\color{refgray}数据：25个州加哥伦比亚特区联合起诉Medicaid工作要求条款\par}
{\small\color{refgray}边际变化：中期选举前不会出现重大政策变动，真正的政策博弈将在选举后展开，2027-2028年为关键政策变化窗口。\par}
\paragraph{DB} 过去一年美联储美元指数几乎未变，外汇波动率处于多年低位，但支撑美元的资金结构正在发生根本性替换：外国官方对美债需求持续走弱，外国持有美债比例从峰值50\%降至约30\%；而外国私人资本对美国公司的涌入创下历史纪录，截至2026年3月美国录得超过6000亿美元净股权流入，净股权流入超过政府和机构债券流入的差额也创下纪录。美元正从“地缘政治锚定”转向“科技周期锚定”，AI和资金基础设施创新正以前所未有的力度将全球股权资本吸入美国。
{\small\color{refgray}数据：外国持有美债比例从峰值50\%降至约30\%\par}
{\small\color{refgray}数据：截至2026年3月美国净股权流入超6000亿美元\par}
{\small\color{refgray}数据：净股权流入超过债券流入的差额创历史纪录\par}
{\small\color{refgray}数据：DTCC本月已开始代币化其持有的超过100万亿美元托管资产\par}
{\small\color{refgray}边际变化：美元强势来源从传统美债避险需求转向对AI主题和美国公司回报的追逐，资金结构改变可能使美元在不确定性事件中不再像过去那样“安全”。\par}
\paragraph{GS} 中国二季度实际GDP环比年化增速从一季度的5.3\%降至3.6\%，零售额和固定资产投资明显走弱。减速中约40\%可归因于财政脉冲转负，30\%来自伊朗战争导致化工和成品油产量下降，其余来自天气和春节消费透支。6月主要活动指标中仅出口、发电量和工业增加值同比增速超5\%，房地产销售和汽车销量同比分别下降14\%和19\%。上半年3D打印机产量同比+48.5\%，锂电池+39.3\%，工业机器人+28.0\%，远高于整体工业增加值5.4\%的增速。政策宽松窗口正在打开，但决策者在安全、科技竞争和财政约束间面临取舍。
{\small\color{refgray}数据：Q2实际GDP环比年化增速从Q1的5.3\%降至3.6\%\par}
{\small\color{refgray}数据：财政脉冲转负贡献约40\%的减速\par}
{\small\color{refgray}数据：6月房地产销售同比-14\%，汽车销量同比-19\%\par}
{\small\color{refgray}数据：上半年3D打印机产量+48.5\%，锂电池+39.3\%，工业机器人+28.0\%\par}
{\small\color{refgray}边际变化：宏观环境从“K型”走向更极端的“K型”，高增长部门集中在AI和高端制造，更广泛的国内需求部门仍在调整。\par}
\paragraph{GS} GE Aerospace 2Q26全面超预期：调整后EPS 2.02美元（共识1.86美元），收入126亿美元（同比+24\%），自由现金流30.3亿美元（共识18.2亿美元）。公司上调全年指引：调整后EPS从7.10-7.40美元上调至7.65-7.85美元，自由现金流从80-84亿美元上调至89-92亿美元。商用发动机与服务收入97亿美元（同比+27\%），售后市场是主要驱动力，反映全球机队利用率高位运行和老旧发动机维修需求累积释放。CES部门营业利润率27.3\%，高于模型预期的26.7\%。
{\small\color{refgray}数据：2Q26调整后EPS 2.02美元（共识1.86美元）\par}
{\small\color{refgray}数据：收入126亿美元，同比+24\%\par}
{\small\color{refgray}数据：自由现金流30.3亿美元（共识18.2亿美元）\par}
{\small\color{refgray}数据：CES收入97亿美元，同比+27\%\par}
{\small\color{refgray}边际变化：自由现金流几乎是市场共识的两倍，公司上调全年指引，航空发动机售后市场结构性需求仍在加速。\par}
\paragraph{GS} 日本财务大臣和首相近期表态考虑引导国内储蓄更多投向本国资产。GS认为两项潜在措施（调整GPIF资产配置、将国债纳入NISA免税账户）可能增加对日本国债的需求，但大部分影响将通过信号效应和正反馈循环实现。GPIF管理约2.8万亿美元资产，当前国内债券配置目标25\%，实际已达27\%，在现行策略下增持理论空间约750亿美元。但这一规模在GPIF日常再平衡中不算大，近期再平衡规模已超过当前可增持上限。真正决定日本利率和波动率的仍是财政与货币政策组合。
{\small\color{refgray}数据：GPIF管理约2.8万亿美元资产\par}
{\small\color{refgray}数据：国内债券配置目标25\%，实际已达27\%\par}
{\small\color{refgray}数据：增持理论空间约750亿美元\par}
{\small\color{refgray}数据：GPIF持有约5\%的日本国债市场\par}
{\small\color{refgray}边际变化：市场对GPIF话题的敏感反应更多在押注信号效应而非实际增持量，储蓄回流并非万能，宏观政策才是关键。\par}
\paragraph{GS} 7月中旬亚洲路演显示，对日本整车厂商和轮胎制造商的关注度正在上升，零部件板块情绪依然偏弱。本田吸引力来自美国高敞口和摩托车业务持续增长；丰田方面，读者问题转向本财年盈利超预期幅度和未来资本配置策略。零部件板块面临中国供应商竞争，本田和三菱等整车厂接连宣布将单车成本削减数十万日元。轮胎板块讨论基调积极，横滨橡胶关注度最高，其市场份额扩张和低成本生产能力是主要吸引力。
{\small\color{refgray}数据：本田美国业务占比高，摩托车销售持续强劲\par}
{\small\color{refgray}数据：本田和三菱等整车厂宣布单车成本削减数十万日元\par}
{\small\color{refgray}数据：横滨橡胶关注度最高，受益于市场份额扩张和低成本生产\par}
{\small\color{refgray}数据：三菱宣布2027年起月产1000台人形机器人\par}
{\small\color{refgray}边际变化：整车板块关注度回升，但零部件板块面临中国供应商全面竞争的压力上升，轮胎板块通过提价稳住局面。\par}
\paragraph{MS} MS提出分析框架：所有生命与非生命系统本质上是将能量（e）转化为智能（i）的过程，公式为 i = er²，其中r代表机器人。对于SpaceX，报告拆解为三个嵌套效率指标：智能每瓦特（i/w）、智能每瓦特每美元（(i/w)/\$）、智能每瓦特每美元每秒（((i/w)/\$)/t）。SpaceX优势来自更高效的芯片设计、垂直整合、数据闭环和制造递归循环。300美元目标价拆分为：空间业务8美元、星链连接128美元、X+Grok 12美元、企业AI 152美元，AI相关占过半估值。
{\small\color{refgray}数据：SpaceX目标价300美元，AI相关占过半（企业AI 152美元，X+Grok 12美元）\par}
{\small\color{refgray}数据：星链连接估值128美元\par}
{\small\color{refgray}数据：空间业务估值8美元\par}
{\small\color{refgray}数据：企业AI估值打了50\%的折扣\par}
{\small\color{refgray}边际变化：MS用能量-智能转化框架重新定义SpaceX价值，强调其竞争优势在于智能密度而非发射次数或卫星数量。\par}
\paragraph{MS} MS策略团队认为美国市场领导力正从半导体和大型科技股向更广泛行业扩散，且尚未结束。支撑因素包括：标普500等权重指数持续跑赢市值加权指数；中位数公司盈利增速出现疫情以来最显著加速；半导体板块盈利修正广度已从历史极端高位回落；油价预期走低缓解通胀担忧；美联储年内按兵不动。半导体板块自6月初已回调超20\%，MS认为修正尚未完成，参照银矿股走势，短期内可能还有约15\%的下行空间。超大规模云厂商相对价值优于半导体，但过去3周已上涨近30\%，性价比不如之前。
{\small\color{refgray}数据：半导体板块自6月初已回调超20\%\par}
{\small\color{refgray}数据：MS认为半导体短期内可能还有约15\%的下行空间\par}
{\small\color{refgray}数据：消费可选和交通运输板块过去两个月跑赢大盘12\%\par}
{\small\color{refgray}数据：超大规模云厂商相对半导体过去3周上涨近30\%\par}
{\small\color{refgray}边际变化：半导体板块的调整是扩散的催化剂而非终结，市场正转向“高质量”方向，但这一过程可能需要几个月时间。\par}
\subsubsection*{关键数据}
\begin{itemize}
  \item 日本科技板块1H26等权重+14\%，市值加权+60\%
  \item 美国软饮市场Q3至今品类增长3.3\%，可口可乐+5.2\%，百事-1.1\%
  \item 外国持有美债比例从峰值50\%降至约30\%，美国净股权流入超6000亿美元
  \item 中国Q2 GDP环比年化增速从5.3\%降至3.6\%，财政脉冲转负贡献40\%
  \item GE Aerospace 2Q26自由现金流30.3亿美元（共识18.2亿美元）
  \item GPIF增持日债理论空间约750亿美元
  \item 半导体板块自6月初已回调超20\%，MS认为可能还有约15\%下行空间
  \item OBBBA法案预计2025-2034年削减约8860亿美元Medicaid支出
\end{itemize}
\begin{figure}[htbp]
\centering
\includegraphics[width=0.88\linewidth]{figures/fig_029.jpg}
\caption{Exhibit 2 - EXHIBIT 2: In 1H26, the sector was up +14\% on equal-weighted basis and +60\% on cap-weighted basis vs. the market, up +15\% Data til Jun 30 \$\textasciicircum{}\{th\}\$ 2026. Returns are in USD}
\end{figure}
\begin{figure}[htbp]
\centering
\includegraphics[width=0.88\linewidth]{figures/fig_067.jpg}
\caption{EXHIBIT 4 - EXHIBIT 4: KO Overall Performance Trends EXHIBIT 5: KO Performance Trends by Segment}
\end{figure}
\begin{figure}[htbp]
\centering
\includegraphics[width=0.88\linewidth]{figures/fig_101.jpg}
\caption{Exhibit 1 - Exhibit 1: Fiscal impulse turned negative in Q2... Exhibit 2: ...Contributing to slower Q2 GDP growth \#\# From Divergent to Very Divergent}
\end{figure}
\begin{figure}[htbp]
\centering
\includegraphics[width=0.88\linewidth]{figures/fig_157.jpg}
\caption{Exhibit 1 - Exhibit 2: Semis Earnings Revisions Breadth Beginning to Decelerate from Upside Extremes.....}
\end{figure}
{\small\color{refgray}本节综合 11 篇研究；涉及 Bernstein、DB、GS、MS。}
\newpage
\section{辅助信号｜战略咨询}
本板块整合波士顿咨询两份报告，聚焦AI结构性成本削减与国防科技利润结构分析，揭示增长与利润的错配及长期竞争格局。
\subsection{AI驱动的结构性成本优势正在拉开企业绩效差距}
\textbf{AI领先者与追赶者之间的绩效差距已结构化，领先企业成本降幅是同行三倍，EBIT利润率高出1.6倍，投入资本回报率高出2.7倍；关键在于AI嵌入核心运营的深度，而非技术选型。}
\begin{itemize}
  \item 仅5\%的公司规模化通过AI创造价值，60\%的公司投入巨大却未获实质性收益；领先者与追赶者差距已反映在利润表中。
  \item AI应用分四个递进阶段：聊天机器人/副驾驶、AI驱动工作流、自主代理、多代理系统；多数公司停留前两阶段，成本优势在第三、四阶段才真正显现。
  \item 五个常见陷阱：碎片化试点无法规模化、AI嫁接在低效流程上、长期项目缺乏早期验证、因等待确定性而停滞、只设效率目标不设成本目标。
  \item 案例：全球消费品公司通过AI部署采购到付款全流程，100天内完成3500次供应商谈判，实现9亿美元总成本削减；全球制药公司应用AI将筛选化合物库扩大100倍，发现周期缩短40\%-50\%。
  \item 边际变化：成本绩效已成为AI成熟度的函数，先发者节省的成本持续投入下一轮能力建设，抬高竞争门槛，追赶者成本将越来越高。
\end{itemize}
\subsection{国防科技新赛道增速惊人但利润池结构决定传统平台长期主导}
\textbf{可消耗系统年增速35\%-40\%，但基数极小，到2033年精尖系统仍占市场超80\%份额；传统平台利润来自数十年持续维护，可持续性远优于新玩家。}
\begin{itemize}
  \item 2025年美国与欧盟采购收入：精尖系统约650亿美元，可负担量产系统约50亿美元，可消耗系统仅5500万美元；到2033年精尖系统年增长2\%-3\%，可负担量产系统15\%-20\%，可消耗系统35\%-40\%。
  \item 精尖系统到2033年仍占市场超80\%份额，新赛道高速增长无法在绝对规模上撼动传统市场，是“补充”而非“颠覆”。
  \item 利润结构差异：精尖系统一半利润来自后期维护（如F/A-18E/F全生命周期利润超10亿美元），而可消耗系统利润主要来自一次性销售，缺乏持续服务收入。
  \item 传统平台利润的“形状”是数十年持续服务，新玩家利润可持续性远低于传统巨头；增长率不等于利润池大小。
  \item 边际变化：国防预算战后大幅波动（美国冷战后下降约37\%，后续下降30\%-45\%），但传统平台维护收入相对稳定，提供利润韧性。
\end{itemize}
\begin{figure}[htbp]
\centering
\includegraphics[width=0.88\linewidth]{figures/fig_180.jpg}
\caption{EXHIBIT 2 - The durability of profits from sustainment for traditional primes is significant. Defense spending sees large swings based on wartime demand. In the US, for example, defense budgets fell roughly 37\% after the Cold War and 30\% to 45\% in the drawdowns following }
\end{figure}
\newpage
\section{辅助信号｜智库/国际机构}
本日智库研究聚焦三大主题：东南亚农业气候与性别双重约束、稳定币驱动的新型美元化、以及全球利润集中与企业税基侵蚀。ADB、BIS、OECD、IMF的报告均指向结构性风险与政策缺口，对新兴市场资产叙事、企业税负预期及金融科技监管具有信号意义。
\subsection{东南亚农业：气候冲击与性别不平等的恶性循环}
\textbf{ADBI指出，气候脆弱性与性别不平等在东南亚小农农业中相互强化，家庭内部资源错配导致女性管理土地产量低约30\%，标准干预因忽视这一机制而失效。}
\begin{itemize}
  \item 东南亚是全球气候不确定性最高地区之一，缅甸、菲律宾、越南位列过去十年极端天气影响最严重十国；该地区约75\%GDP依赖自然资本相关部门。
  \item 同一家庭内女性管理土地产量比男性低约30\%，原因并非女性效率低，而是家庭将过多劳动力和过少肥料分配给男性地块，错配造成家庭农业总产出约6\%的效率损失。
  \item 印度尼西亚仅13\%女性拥有或享有农业土地稳定权利，男性为52\%；气候冲击时资源受限家庭更倾向将投入集中于男性土地，进一步扩大生产力差距。
  \item 传统发展项目多基于“统一家庭模型”设计，但过去三十年实证研究已否定该假设，导致干预措施效果不彰。
\end{itemize}
\begin{figure}[htbp]
\centering
\includegraphics[width=0.88\linewidth]{figures/fig_167.jpg}
\caption{Figure 1 - Figure 1: Climate Vulnerability, Gender Inequality, and Maladaptation in Southeast Asian Agriculture Creates a Self-Reinforcing Hotspot \#\# QUANTIFYING THE DUAL CONSTRAINT \#\# Climate Vulnerability and Agriculture}
\end{figure}
\subsection{稳定币与新型美元化：对资本管制免疫且难以逆转}
\textbf{BIS工作论文系统比较传统存款美元化与稳定币美元化，发现两者驱动因素高度相似，但稳定币对资本管制几乎免疫且一旦形成即高度持久，对新兴市场货币政策构成新挑战。}
\begin{itemize}
  \item 基于130个地区1990-2019年存款美元化数据及2017-2024年跨境稳定币流入数据，两者均与汇率传递效应强度、主权或银行扰动显著相关。
  \item 银行扰动对稳定币流入的推动作用尤为突出，符合稳定币的非银行属性——银行体系不稳定时，非银行美元获取渠道更具吸引力。
  \item 两种美元化均高度持久：存款美元化具有显著自回归特征，稳定币初步证据同样显示高度持续性，意味着美元化一旦建立很难逆转。
  \item 存款美元化与稳定币之间替代效应有限，暗示两个市场存在一定割裂，稳定币可能成为银行体系脆弱性的放大器。
\end{itemize}
\subsection{全球利润集中与企业税基侵蚀：研究中心辖区人均利润显著偏高}
\textbf{OECD 2026年企业税统计显示，全球大型跨国企业利润仍高度集中于研究中心辖区，关联方收入占比偏高，提示税基侵蚀与利润转移行为持续存在。}
\begin{itemize}
  \item 基于约9400家跨国企业集团国别报告数据，研究中心辖区（内向FDI存量超GDP150\%）员工人均收入中位数181.1万美元，高收入辖区47.7万，中等收入21.1万，低收入仅15.3万。
  \item 研究中心辖区关联方收入占总收入比例中位数超30\%，高收入辖区20\%，中等收入14\%，低收入7\%；报告明确表示高水平关联方收入是税务筹划的高不确定性评估因素。
  \item 人均收入与关联方收入占比两个指标结合观察，比单独看税率或利润总额更有意义，研究中心高人均收入可能来自利润集中配置而非更高运营效率。
  \item 全球最低税实施前，不同收入组别企业税负差异显著，低收入国家企业税占GDP可低至1\%以下，而投资枢纽地区税负比例反而更高。
\end{itemize}
\begin{figure}[htbp]
\centering
\includegraphics[width=0.88\linewidth]{figures/fig_175.jpg}
\caption{Figure 1 - \textbackslash{}- the current stage of the economic cycle and the degree of cyclicality of the corporate tax system (for example, from the generosity of loss offset provisions); \textbackslash{}- other instruments that postpone the taxation of earned profits. Generally, differences in corp}
\end{figure}
\subsection{AI、技术与生产率}
\textbf{用于判断 AI 投资周期、生产率扩散和产业约束是否正在改变资产定价。}
\begin{itemize}
  \item 亚洲开发银行（ADB）于2026年7月发布《社会保护覆盖与效果指标：社会保护评估工具》，提出了一套名为SPICES的标准化评估框架。该框架的核心变化在于，将社会保护评估的焦点从“投入了多少资金”转向“这些资金覆盖了哪些人群、解决了哪些需求、产生了什么效果”。报告认为，亚太地区各国...
  \item 2025年，经合组织成员国中超过三分之一的成年人报告使用AI工具。这些工具正越来越多地被用于支持资金决策——从选择资金产品、预算管理、信贷管理，到研究和退休规划。消费者还依赖AI获取个性化的资金教育，有时会提出他们不习惯向人类顾问询问的问题。
\end{itemize}
\begin{figure}[htbp]
\centering
\includegraphics[width=0.88\linewidth]{figures/fig_169.jpg}
\caption{Figure 1 - Figure 2: Areas of Need Requiring Social Protection}
\end{figure}
\subsection{宏观政策与金融稳定}
\textbf{用于校准利率、通胀、财政约束和金融稳定叙事。}
\begin{itemize}
  \item 2024年，冈比亚人均收入约为940美元，仍不足撒哈拉以南非洲地区平均水平的一半。这份由IMF工作人员完成的国别报告指出，该国过去三十年的宏观环境增长主要依靠劳动力和资本积累驱动，全要素生产率的贡献波动且时常为负。随着人口结构变化和资本积累空间收窄，冈比亚需要加速结构性改革以提升...
\end{itemize}
\begin{figure}[htbp]
\centering
\includegraphics[width=0.88\linewidth]{figures/fig_172.jpg}
\caption{Figure 1 - 3. This paper highlights past, current and future drivers of The Gambia's growth and suggests reforms to enhance productivity. Section B documents recent economic developments. A growth accounting is undertaken in Section C to identify the contributions of fac}
\end{figure}
\section{结语}
后续需验证的关键节点包括：美联储7月会议对利率路径的暗示、美国科技股财报对AI资本开支的指引、中国7月政治局会议对财政与地产政策的定调、卡塔尔LNG冬季供应恢复的实际进展、以及美国中期选举前医疗政策博弈的演变。
\newpage
\section{覆盖概览}
投行/券商 55 篇；战略咨询 2 篇；智库/国际机构 6 篇。\par
投行覆盖：Bernstein 20篇 / GS 14篇 / MS 7篇 / NOM 6篇 / DB 3篇 / JPM 3篇 / Citi 1篇 / UBS 1篇\par
\section{Disclaimer}
{\scriptsize\color{refgray}
This community is a paid learning and information-sharing community created for general educational purposes only. The membership fee is charged solely for access to community discussions, educational content, organization, and operational costs. It is not an advisory fee, management fee, performance fee, brokerage fee, or compensation for personalized financial, investment, legal, tax, or professional advice.\par
I participate and share information solely as an individual and for educational discussion among community members. I am not acting as a financial adviser, investment adviser, broker, dealer, portfolio manager, legal adviser, tax adviser, accountant, fiduciary, or any other licensed professional adviser. Nothing shared in this community constitutes, and should not be interpreted as, financial advice, investment advice, legal advice, tax advice, a recommendation, solicitation, offer, endorsement, instruction, or guarantee to buy, sell, hold, short, trade, or invest in any security, cryptocurrency, fund, derivative, strategy, or other financial product.\par
All content shared in this community, including messages, discussions, links, files, charts, examples, market commentary, watchlists, AI-generated or AI-polished summaries, and third-party materials, is general, impersonal, and educational in nature. It does not take into account any individual's financial situation, investment objectives, risk tolerance, investment horizon, liquidity needs, tax status, legal restrictions, or suitability requirements. No content should be relied upon as suitable for any specific person, account, portfolio, or financial decision.\par
Information shared in this community may be incomplete, inaccurate, outdated, speculative, AI-generated, AI-edited, or based on third-party internet sources that have not been independently verified. I make no representation or warranty as to the accuracy, completeness, timeliness, reliability, or suitability of any information shared.\par
Financial markets involve substantial risk, including the possible loss of principal. Past performance, historical data, hypothetical examples, simulated results, backtests, personal opinions, or educational examples do not guarantee future results. Each member is solely responsible for conducting their own research, exercising independent judgment, and making their own decisions. Before making any financial, investment, legal, or tax decision, members should consult qualified and licensed professionals.\par
Participation in this community, payment of any membership fee, or communication with me or other members does not create any adviser-client, fiduciary, professional, contractual, agency, partnership, employment, or other advisory relationship. By joining or participating in this community, you acknowledge and agree that you are solely responsible for your own decisions, actions, transactions, profits, losses, risks, and consequences, and that I shall not be liable for any direct, indirect, incidental, consequential, financial, legal, tax, or other loss or damage arising from or related to any content, discussion, information, or material shared in this community.\par
}
\newpage
\begin{center}
{\Large 更多详情报告kcdesk.com}\par\vspace{0.8cm}
\includegraphics[width=0.58\linewidth]{\detokenize{../../prompts/zsxq_img.jpg}}
\end{center}
\end{document}
